% ─── Chapter 6: Training Infrastructure & Experiment Tracking ─────
\chapter{Training Infrastructure \& Experiment Tracking}
\label{ch:infrastructure}

\section{Distributed Training Architecture}

Training is orchestrated by Ray RLlib, which distributes rollout collection
across multiple workers and performs SGD updates on a central learner:

\begin{center}
\small
\begin{tabularx}{\textwidth}{l X}
  \toprule
  \textbf{Component} & \textbf{Role} \\
  \midrule
  Main process   & Configures training, manages curriculum callbacks,
                   drives the training loop. \\
  GPU learner     & Performs PPO SGD updates on collected batches.
                   Uses PyTorch backend. \\
  Rollout workers & Each runs multiple environments, collects trajectories,
                   and sends batches to the learner. \\
  Showdown servers & 8 local Node.js instances on ports 8000--8007,
                     each hosting parallel battles. \\
  \bottomrule
\end{tabularx}
\end{center}

The default (\texttt{standard}) configuration uses 12~workers with 4
environments per worker, yielding \textbf{48~parallel battles}.  Each Showdown
server hosts approximately 6~concurrent battles.

\section{Hardware}

\begin{center}
\small
\begin{tabularx}{0.8\textwidth}{l l}
  \toprule
  \textbf{Resource} & \textbf{Specification} \\
  \midrule
  GPU         & NVIDIA RTX 5090 (\texttt{optimal} preset) \\
  RAM         & $\sim$56\,GB peak with 48 environments \\
  CPU         & Multi-core (workers run on CPU threads) \\
  Storage     & Local SSD for checkpoints and MLflow artifacts \\
  \bottomrule
\end{tabularx}
\end{center}

Memory management is critical: the project previously observed RAM growth from
8\,GB to 52\,GB over 5~hours due to stale entries in environment caches.  This
was mitigated with periodic cleanup and memory-leak monitoring.

\section{Configuration Presets}

Five hardware presets control model size, parallelism, and training duration:

\begin{center}
\small
\begin{tabularx}{\textwidth}{l c c c c}
  \toprule
  \textbf{Parameter} & \textbf{Quick} & \textbf{Standard} & \textbf{Memory Safe} & \textbf{Large} \\
  \midrule
  Hidden dim              & 128   & 512   & 256   & 768   \\
  Attention heads         & 8     & 8     & 4     & 12    \\
  Transformer layers      & 1     & 1     & 1     & 6     \\
  Embedding dim           & 32    & 32    & 32    & 32    \\
  LSTM hidden             & ---   & 512   & 256   & 768   \\
  Use LSTM                & No    & Yes   & No    & Yes   \\
  Train batch size        & 4\,096  & 6\,144  & 2\,048  & 16\,384 \\
  SGD minibatch size      & 256   & 256   & 128   & 512   \\
  Workers                 & 0     & 12    & 4     & 12    \\
  Env/worker              & 1     & 4     & 2     & 4     \\
  Total envs              & 1     & 48    & 8     & 48    \\
  Total timesteps         & 150K  & 10M   & 10M   & 10M   \\
  Approx.\ params         & $\sim$1M & $\sim$10M & $\sim$4M & $\sim$42M \\
  \bottomrule
\end{tabularx}
\end{center}

The \texttt{quick} preset is used for smoke tests and debugging.  The
\texttt{standard} preset is the default for training runs.  The
\texttt{optimal} preset (not shown) matches \texttt{standard} but targets the
RTX~5090 specifically.

\section{PPO Hyperparameters}

Default hyperparameters across all presets (unless overridden):

\begin{center}
\small
\begin{tabularx}{0.7\textwidth}{l r}
  \toprule
  \textbf{Parameter} & \textbf{Value} \\
  \midrule
  Learning rate              & $3 \times 10^{-4}$ \\
  Discount ($\gamma$)        & 0.95 \\
  GAE lambda ($\lambda$)     & 0.95 \\
  Clip range ($\epsilon$)    & 0.2 \\
  Entropy coefficient        & 0.05 \\
  Value-function loss coeff. & 0.5 \\
  Value clip param           & 5.0 \\
  Gradient clipping          & 5.0 \\
  SGD iterations per batch   & 10 \\
  \bottomrule
\end{tabularx}
\end{center}

\section{MLflow Experiment Tracking}

All training runs are logged to an \textbf{MLflow} tracking server under the
experiment name \texttt{Pokemon\_RL\_Battler}.  The following categories of
data are recorded:

\begin{itemize}
  \item \textbf{Configuration snapshot:} Full training config, model config,
        curriculum config, and hardware preset---stored as MLflow parameters.
  \item \textbf{Training metrics:} Policy loss, value loss, entropy,
        KL divergence, win rate, episode length, per-stage statistics---logged
        every training iteration.
  \item \textbf{Validation metrics:} Win rates against fixed opponents
        (random, heuristic), per-opponent-type breakdowns, fallback event
        counts---logged at scheduled evaluation intervals.
  \item \textbf{Diagnostic metrics:} Action-mask density, observation sample
        recordings, HP remaining distributions, faint counts per episode.
  \item \textbf{System metrics:} CPU/RAM/GPU utilisation, memory leak counters.
\end{itemize}

Metrics use nested keys (e.g.\
\texttt{validation/fixed\_paired/win\_rate\_vs\_heuristic}) to support
structured comparison in the MLflow UI.

% TODO: Insert MLflow screenshot showing the experiment comparison view
% with multiple runs.  Highlight win rate curves, loss curves, and
% the ability to compare configuration differences between runs.

\section{Checkpoint Management}

\begin{itemize}
  \item \textbf{Frequency:} Checkpoints saved every 500\,000 training steps.
  \item \textbf{Retention:} Up to 5~checkpoints retained (oldest automatically
        deleted).
  \item \textbf{Self-play export:} At each checkpoint, model weights are also
        exported to \texttt{checkpoints/selfplay\_latest.pt} for use by
        \texttt{SelfPlayPlayer} instances.
  \item \textbf{Resume:} Training can be resumed from the latest checkpoint
        using \texttt{-{}-resume-checkpoint latest}, preserving MLflow run
        continuity via \texttt{-{}-mlflow-run-id <RUN\_ID>}.
\end{itemize}

\section{Validation Pipeline}

Scheduled evaluation runs against fixed opponent configurations:

\begin{center}
\small
\begin{tabularx}{\textwidth}{l X}
  \toprule
  \textbf{Protocol} & \textbf{Description} \\
  \midrule
  Smoke test     & Quick sanity check: 10 episodes against random opponent. \\
  Fixed paired   & 20~frozen teams $\times$ 2~opponents (random + heuristic) =
                   40~battles per evaluation run.  Primary comparison metric. \\
  Mirror         & Both sides use the same team; tests raw battling skill
                   without team advantage. \\
  \bottomrule
\end{tabularx}
\end{center}

Validation is triggered at regular intervals (every 100\,000 steps) during
training and can also be run manually against any checkpoint via
\texttt{scripts/validate\_checkpoint.py}.  Results are logged to MLflow with
per-opponent-type breakdowns.

\section{Training Throughput}

Observed throughput on the \texttt{standard} preset:
$\sim$800~environment samples per second (target: 1\,000+).  The primary
bottleneck is action-mask computation, which accounts for approximately 82\,\%
of per-step wall time (2.3\,ms masking vs.\ 0.5\,ms forward pass).
